\chapter*{Introduction Générale}
\addcontentsline{toc}{chapter}{Introduction Générale}

Les centres opérationnels de sécurité reçoivent quotidiennement des informations décrivant des campagnes malveillantes, des infrastructures adverses, des fichiers suspects ou des modes opératoires. Ces données de renseignement sur les menaces, désignées par le terme \textit{Cyber Threat Intelligence} (CTI), ne deviennent toutefois utiles à la surveillance que lorsqu'elles sont transformées en mécanismes de détection adaptés aux journaux réellement disponibles. Cette transformation demande habituellement plusieurs opérations : normalisation des événements, extraction des comportements, association aux techniques MITRE ATT\&CK, vérification de la télémétrie, recherche de doublons, rédaction de règles Sigma, compilation, validation et publication.

Lorsque ces opérations sont réalisées manuellement, leur enchaînement peut être long et difficile à reproduire. La qualité dépend en outre de la disponibilité de l'analyste et de sa connaissance simultanée du renseignement, des sources de journaux et des conventions de détection. L'intelligence artificielle générative peut accélérer certaines tâches de synthèse et de rédaction, mais ses résultats ne doivent pas être considérés comme fiables sans contrôle. Une règle plausible sur le plan linguistique peut contenir un champ inexistant, une association ATT\&CK incorrecte, une logique trop large ou une information absente de l'événement source.

Le stage d'un mois effectué au sein de Forvis Mazars Group s'inscrit dans ce contexte. Il porte sur une plateforme locale d'ingénierie de détection CTI assistée par intelligence artificielle. La solution reçoit de véritables événements MISP, les normalise, orchestre leur traitement dans un graphe LangGraph, produit des propositions Sigma et les soumet à une série de contrôles déterministes. Elle fournit ensuite à l'analyste une file de revue permettant d'approuver, de rejeter ou de demander une correction. Le principe directeur du projet est le suivant : l'intelligence artificielle peut proposer, réparer et recommander, mais elle ne peut ni approuver ni déployer une détection. La décision finale demeure humaine.

L'objectif de ce rapport est de présenter le besoin, les choix de conception, la réalisation et les résultats obtenus. Le premier chapitre expose le cadre du stage, la problématique et les objectifs fonctionnels et techniques. Le deuxième chapitre décrit les notions nécessaires, l'architecture de la plateforme, le modèle de données et le workflow de raisonnement. Le troisième chapitre détaille la mise en œuvre, les mécanismes de validation, l'intégration DevSecOps, les tests, les difficultés rencontrées ainsi que les limites de la solution. Une conclusion générale synthétise les apports du stage et propose des perspectives d'amélioration.